% !TeX root = ../papers/nc-p3.tex
\section{単純点を持つ一辺の十九点模型と全延長の線型核}
\label{sec:nineteen-point-edge}
二つの三次元空間を用い、$X=\mathbb P(V^*)\times\mathbb P(U)$ の点を行 $l$ と列 $z$ で表す。
$T_1,T_2$ を両平面の接束とする。
\begin{theorem}\label{thm:nineteen-point-edge}
左右四文字・六二次関係の一辺の点対が有限で、単純な一点を指定する。
その印を保つ一辺の基底変更類は、二平面の積内の長さ19のスキーム $Y$ から復元される。
現れる $Y$ は
\[
 E_1=T_1\boxtimes\mathcal O(1),\qquad E_2=\mathcal O(1)\boxtimes T_2
\]
の切断を一つずつ取った同時零スキームで、零次元であるものにちょうど一致する。
非被約な $Y$ を許す。二つの切断は $Y$ からそれぞれ非零倍率まで一意である。
\end{theorem}
\begin{proof}
印をそれぞれの星座標点とする。単純点なので局所アフィン六変数について六双一次式の
一次部行列が可逆である。関係の基底を取り直し、座標を $(\gamma,l),(\alpha,z)$ と書けば六式は
\[
 l(\alpha I+F_z)=0,\qquad (\gamma I+Q_l)z=0
\]
となる。ここで $F_z$ と $Q_l$ は添字に線型な三次行列である。
$\alpha,\gamma$ に三座標の一次式を加えることで各々跡零とする。
Euler列によればこの二つの跡零行列系がちょうど上の二束の大域切断である。
星点以外では $l,z$ とも零でない。片方でも零ならもう片方も零となるからである。
二平面に射影した像は二つの固有ベクトル条件であり、固有値 $-\alpha,-\gamma$ は
各単元座標で割ればスキーム上でも一意に戻る。逆に任意の二切断からこの式を作ると星点は単純である。
星点の補集合が有限なら零スキームは階数4の束の正則切断となる。
二平面の超平面類を $h,k$ とすれば最高Chern類は
$(3h^2+3hk+k^2)(h^2+3hk+3k^2)=19h^2k^2$ である。

より強く $h^0(\mathcal I_Y E_i)=1$ を示す。$E=E_1\oplus E_2$ の正則切断のKoszul解消を $E_1$ でひねる。
$H^0(E_1\otimes E^*)$ は恒等写像だけである。実際平面接束の自己準同型はスカラーで、
異なる束間には $H^0(\Omega^1(1))=0$ を使える。
さらに $j=2,3,4$ で $H^{j-1}(E_1\otimes\bigwedge^j E^*)=0$ である。
$j=2$ の三項は順に
$T(-3)\boxtimes\mathcal O(-1)$、$\operatorname{End}(T)(-1)\boxtimes\Omega^1$、
$T(-2)\boxtimes\mathcal O(-2)$ となる。
$j=3$ では $T(-4)\boxtimes\Omega^1(-1)$ と
$\operatorname{End}(T)(-2)\boxtimes\mathcal O(-3)$、
$j=4$ では $T(-5)\boxtimes\mathcal O(-4)$ である。
Euler列とKünneth公式で所要の消滅が従う。
最後は三個の一次式による $H^0(\mathcal O(1))^3\to H^0(\mathcal O(2))$ の全射性を双対化する。
また $H^0(\operatorname{End}T(-1))=H^0(\operatorname{End}T(-2))=0$ は
Euler列、または接束の安定性から従う。よってイデアルに値を持つ切断は元の切断の倍数だけで、
二束を交換しても同じである。二つの倍率は星座標の独立な倍率で吸収される。
印を保つ射影変更は二平面で射影変更を誘導するので、$Y$ の順序付き二因子の射影同値が一辺の同型の必要十分条件である。
\end{proof}

\begin{theorem}\label{thm:nineteen-point-linear-companion}
この一辺を位置 $1,2$ に置く。両反対の印の縮約が階数3の四周期の印付き二切断補完は、
跡零の $F,Q$ を上の十九点から取り、次の線型核で表示できる。
\[
 K_{F,Q}=\{(G,P):[G_a,F_c]_{bd}+[Q_b,P_d]_{ac}=0\}.
\]
非零な組の共通倍率を除き、同じ $Y$ の安定群からの同一視だけが残る。
この表示では有限な標準Frobenius・Koszul左右ASとなるものへの条件を引き続き課す。
特に四辺が有限な標準Frobenius候補に単純な印があればこの表示で尽くされる。
十五次元成分の一般点の全二十印・四始点では $\dim K_{F,Q}=1$ なので、
十九点の配置だけで印付き窓そのものの添字保持同型類が決まる。
\end{theorem}
\begin{proof}
反対二行列を単位行列とし、星文字の持上げで残る二組も跡零にできる。
定理\ref{thm:marked-point-commutators}の二切断のSchur補の計算は、既に分かった $F,Q$ に
ついて読めば、48未知係数 $(G,P)$ の81本の一次式である。
解ごとに星のない四成分を同定理の積和で置けば、両切断がちょうど階数6のテンソルに一意に戻る。
$G=P=0$ の補完も式上は存在するが、反対の辺では星座標を両方とも零とした二平面の積が点対に入り四辺有限の候補から外れる。
非零解では、$F,Q$ を固定したまま位置0,3の三文字を $h$ 倍し、全体を $h^{-1}$ 倍すれば
$(G,P)\mapsto(hG,hP)$ となる。このほかの同一視は元の印付き一辺から誘導される。
具体的には $Y$ を保つ両平面の変換を持ち上げ、二切断の単位行列と $F,Q$ の規格化へ戻して核に作用させる。
持上げの違いは既に除いた共通倍率に過ぎない。
有限Frobenius候補の単純な印で反対縮約がともに階数3となることは
定理\ref{thm:finite-cubic-separators}である。
\path{computations/verify_nineteen_point_edge.py} は分裂混合窓の素数1009の横断二十点を用い、
全八十通りでこの線型写像の階数が47であることを確認する。上界は実際の非零解で従う。
エタールな点の近傍で小行列を持ち上げて一般点での主張を得る。
\end{proof}
